% Slug: ch16-join
% Source: chapters/ch16.tex:172 (tikzpicture env start line)
% Original label: \label{fig:ch16-join}
% Original caption (LaTeX source):
%   The join block recomputes a sum available on every path, so it is fully redundant once the branches reuse the entry value.
% Original caption (plain text):
%   The join block recomputes a sum available on every path, so it is fully redundant once the branches reuse the entry value.
% Required TikZ setup: \usepackage{tikz} + \usetikzlibrary{arrows.meta, positioning, shapes.geometric, fit} (loaded by support/compiler-pseudocode.sty).
% Shared styles used in this picture: `basic block', `cfg edge'.
%   Defined in support/compiler-pseudocode.sty: `basic block', `cfg edge', `back edge', `block label'.
% Usage: % Slug: ch16-join
% Source: chapters/ch16.tex:172 (tikzpicture env start line)
% Original label: \label{fig:ch16-join}
% Original caption (LaTeX source):
%   The join block recomputes a sum available on every path, so it is fully redundant once the branches reuse the entry value.
% Original caption (plain text):
%   The join block recomputes a sum available on every path, so it is fully redundant once the branches reuse the entry value.
% Required TikZ setup: \usepackage{tikz} + \usetikzlibrary{arrows.meta, positioning, shapes.geometric, fit} (loaded by support/compiler-pseudocode.sty).
% Shared styles used in this picture: `basic block', `cfg edge'.
%   Defined in support/compiler-pseudocode.sty: `basic block', `cfg edge', `back edge', `block label'.
% Usage: % Slug: ch16-join
% Source: chapters/ch16.tex:172 (tikzpicture env start line)
% Original label: \label{fig:ch16-join}
% Original caption (LaTeX source):
%   The join block recomputes a sum available on every path, so it is fully redundant once the branches reuse the entry value.
% Original caption (plain text):
%   The join block recomputes a sum available on every path, so it is fully redundant once the branches reuse the entry value.
% Required TikZ setup: \usepackage{tikz} + \usetikzlibrary{arrows.meta, positioning, shapes.geometric, fit} (loaded by support/compiler-pseudocode.sty).
% Shared styles used in this picture: `basic block', `cfg edge'.
%   Defined in support/compiler-pseudocode.sty: `basic block', `cfg edge', `back edge', `block label'.
% Usage: % Slug: ch16-join
% Source: chapters/ch16.tex:172 (tikzpicture env start line)
% Original label: \label{fig:ch16-join}
% Original caption (LaTeX source):
%   The join block recomputes a sum available on every path, so it is fully redundant once the branches reuse the entry value.
% Original caption (plain text):
%   The join block recomputes a sum available on every path, so it is fully redundant once the branches reuse the entry value.
% Required TikZ setup: \usepackage{tikz} + \usetikzlibrary{arrows.meta, positioning, shapes.geometric, fit} (loaded by support/compiler-pseudocode.sty).
% Shared styles used in this picture: `basic block', `cfg edge'.
%   Defined in support/compiler-pseudocode.sty: `basic block', `cfg edge', `back edge', `block label'.
% Usage: \input{graphs-code/ch16-join.tikz.tex} inside a figure environment
%   (wrap with \begin{figure} ... \caption{...} \label{fig:ch16-join} ... \end{figure}).
% Note: tikzpicture body below is verbatim from the source; do not reformat.
\begin{tikzpicture}[node distance=14mm and 18mm]
    \node[basic block] (ch16b0) {t0 = a + b};
    \node[basic block,below left=of ch16b0] (ch16b1) {t1 = a + b};
    \node[basic block,below right=of ch16b0] (ch16b2) {t2 = a + b};
    \node[basic block,below right=of ch16b1] (ch16b3) {t3 = a + b};
    \draw[cfg edge] (ch16b0) -- (ch16b1);
    \draw[cfg edge] (ch16b0) -- (ch16b2);
    \draw[cfg edge] (ch16b1) -- (ch16b3);
    \draw[cfg edge] (ch16b2) -- (ch16b3);
  \end{tikzpicture}
 inside a figure environment
%   (wrap with \begin{figure} ... \caption{...} \label{fig:ch16-join} ... \end{figure}).
% Note: tikzpicture body below is verbatim from the source; do not reformat.
\begin{tikzpicture}[node distance=14mm and 18mm]
    \node[basic block] (ch16b0) {t0 = a + b};
    \node[basic block,below left=of ch16b0] (ch16b1) {t1 = a + b};
    \node[basic block,below right=of ch16b0] (ch16b2) {t2 = a + b};
    \node[basic block,below right=of ch16b1] (ch16b3) {t3 = a + b};
    \draw[cfg edge] (ch16b0) -- (ch16b1);
    \draw[cfg edge] (ch16b0) -- (ch16b2);
    \draw[cfg edge] (ch16b1) -- (ch16b3);
    \draw[cfg edge] (ch16b2) -- (ch16b3);
  \end{tikzpicture}
 inside a figure environment
%   (wrap with \begin{figure} ... \caption{...} \label{fig:ch16-join} ... \end{figure}).
% Note: tikzpicture body below is verbatim from the source; do not reformat.
\begin{tikzpicture}[node distance=14mm and 18mm]
    \node[basic block] (ch16b0) {t0 = a + b};
    \node[basic block,below left=of ch16b0] (ch16b1) {t1 = a + b};
    \node[basic block,below right=of ch16b0] (ch16b2) {t2 = a + b};
    \node[basic block,below right=of ch16b1] (ch16b3) {t3 = a + b};
    \draw[cfg edge] (ch16b0) -- (ch16b1);
    \draw[cfg edge] (ch16b0) -- (ch16b2);
    \draw[cfg edge] (ch16b1) -- (ch16b3);
    \draw[cfg edge] (ch16b2) -- (ch16b3);
  \end{tikzpicture}
 inside a figure environment
%   (wrap with \begin{figure} ... \caption{...} \label{fig:ch16-join} ... \end{figure}).
% Note: tikzpicture body below is verbatim from the source; do not reformat.
\begin{tikzpicture}[node distance=14mm and 18mm]
    \node[basic block] (ch16b0) {t0 = a + b};
    \node[basic block,below left=of ch16b0] (ch16b1) {t1 = a + b};
    \node[basic block,below right=of ch16b0] (ch16b2) {t2 = a + b};
    \node[basic block,below right=of ch16b1] (ch16b3) {t3 = a + b};
    \draw[cfg edge] (ch16b0) -- (ch16b1);
    \draw[cfg edge] (ch16b0) -- (ch16b2);
    \draw[cfg edge] (ch16b1) -- (ch16b3);
    \draw[cfg edge] (ch16b2) -- (ch16b3);
  \end{tikzpicture}
